\chapter{CONCLUSION AND FUTURE WORK}

This chapter summarizes the outcomes of the proposed intelligent 
autoscaling framework for the microservice-based course registration 
system and outlines possible directions for future research. The 
study aimed to design a robust, queuing-theoretically grounded, and 
operationally transparent system capable of dynamically scaling 
auth-service, course-service, and seat-service containers in response 
to fluctuating traffic demands. The approach combined the 
Multi-Feature Decision System (MFDS), Minimum Scaling Overhead (MSO) 
engine, Hysteresis Guard, and Flexible Probability Routing (FPR) 
engine to address scaling oscillation, SLA violations, and resource 
over-provisioning challenges inherent in threshold-based autoscaling 
approaches.

% ─────────────────────────────────────────────────────────────────
\section{\uppercase{Conclusion}}

The developed framework successfully integrated several 
state-of-the-art components, including real-time Prometheus-based 
metrics collection, M/M/N queue-theoretic decision making, 
MSO-driven direct replica computation, hysteresis-based 
anti-oscillation control, and FPR-based proportional traffic 
routing. Among all evaluated configurations, the full 
MFDS-MSO-FPR pipeline with hysteresis demonstrated superior 
performance with an overall SLA violation reduction of 91.7\% 
compared to the baseline reactive scaling approach, achieving 
only 1 SLA violation across all five traffic phases compared to 
12 in the baseline. The results confirm that queuing-theoretic 
autoscaling, when combined with direct scaling computation and 
hysteresis confirmation, maintains high responsiveness and 
stability across all traffic conditions.

The Hysteresis Guard played a vital role in eliminating scaling 
oscillations, requiring three consecutive confirmations of 
under-utilization before executing any scale-in action. This 
mechanism effectively prevented the thrashing behavior observed 
in the baseline system, where 8 oscillations were recorded across 
the five traffic phases. The MSO direct-jump formula further 
ensured that scale-out actions reached the mathematically optimal 
replica count in a single Docker scaling command rather than 
through incremental step-by-step adjustments, significantly 
reducing the time spent in an under-provisioned state during 
traffic spikes. Furthermore, the real-time dashboard provided 
valuable operational transparency, confirming that the system's 
scaling decisions corresponded to the computed queue-theoretic 
parameters at every evaluation cycle.

Overall, this research established a reproducible and 
operationally transparent autoscaling pipeline that aligns with 
cloud-native deployment objectives, ensuring SLA compliance, 
resource efficiency, and stability of microservice workloads 
under dynamic traffic conditions.

% ─────────────────────────────────────────────────────────────────
\section{\uppercase{Key Contributions}}

The primary contributions of this research can be summarized as 
follows. First, it introduced a unified intelligent autoscaling 
system combining M/M/N queuing theory and direct scaling 
computation to overcome the limitations of threshold-based and 
reactive autoscaling approaches in microservice environments. 
Second, the study implemented and evaluated the complete 
MFDS-MSO-FPR pipeline under controlled multi-phase traffic 
conditions, providing a comprehensive understanding of each 
algorithm's contribution to overall scaling performance. Third, 
the work demonstrated that the Hysteresis Guard with a 
three-cycle confirmation threshold effectively eliminates scaling 
oscillations without introducing significant response delay, 
thereby improving the stability and fairness of autoscaling 
decisions. Fourth, the MSO direct-jump formula was shown to 
reduce scale-out lag by computing the optimal replica count 
directly from the M/M/N utilization target, eliminating the 
incremental scaling steps that cause prolonged SLA violations 
during sudden traffic spikes. Finally, the real-time monitoring 
dashboard with Topology, Metrics, and Events tabs provided full 
operational visibility into every queue-theoretic parameter and 
scaling decision, contributing to operator trust and 
transparency in AI-driven infrastructure management.

% ─────────────────────────────────────────────────────────────────
\section{\uppercase{Limitations}}

While the overall results were promising, certain limitations must be acknowledged. The evaluation was conducted in a controlled local Docker Swarm environment with simulated traffic patterns, which may not fully capture the complexity and unpredictability of real-world production traffic at scale. The M/M/c queue model assumes Poisson arrival rates and exponential service times, which are reasonable approximations but may deviate from actual traffic distributions in highly bursty or correlated workloads.

The current system relies primarily on HTTP request rate and latency metrics from Prometheus; integration with CPU, memory, and network I/O metrics would provide a more comprehensive view of service health. The fixed cooldown period of 90 seconds, while effective at preventing oscillation, introduces a predetermined delay that may be suboptimal for workloads with very rapid traffic transitions.\\

Additionally, the framework was evaluated on three microservices with relatively uniform service time characteristics. Real-world microservices architectures typically involve heterogeneous services with significantly varying processing complexities and resource requirements, which would require service-specific queue model calibration and potentially different queuing model types to accurately represent their behavior.

% ─────────────────────────────────────────────────────────────────
\section{\uppercase{Future Work}}

Future research may focus on large-scale multi-service validation 
to ensure the MFDS-MSO-FPR pipeline generalizes across diverse 
microservice architectures with heterogeneous workload 
characteristics. Integrating multimodal metrics such as CPU 
utilization, memory pressure, network throughput, and error 
rates can provide a more comprehensive input feature set for 
the MFDS decision engine, enabling more accurate queue state 
estimation. Emerging machine learning approaches such as 
reinforcement learning-based autoscaling and predictive scaling 
using LSTM or Transformer-based time-series forecasting offer 
potential for proactive replica adjustment before traffic spikes 
occur, rather than reacting after the fact. Exploring adaptive 
SLA bounds that dynamically adjust $T_{min}$ and $T_{max}$ 
based on time-of-day traffic patterns and historical workload 
profiles will further improve resource efficiency. Further 
enhancement of the routing layer by integrating service mesh 
technologies such as Istio or Envoy can enable fine-grained 
per-request traffic shaping beyond the current NGINX upstream 
configuration. Finally, deploying the system on a managed 
Kubernetes cluster with Horizontal Pod Autoscaler (HPA) 
integration can facilitate real-world cloud-native deployment, 
enabling validation against production-grade traffic at scale 
on platforms such as Amazon EKS or Google GKE.\\

% ─────────────────────────────────────────────────────────────────
\section{\uppercase{Closing Remarks}}

In summary, this project successfully demonstrated that combining queuing-theoretic decision making through the MFDS algorithm, direct scaling computation via the MSO strategy, and proportional traffic routing using FPR yields a powerful approach to intelligent autoscaling for microservice-based systems. The integrated framework addresses fundamental limitations of traditional threshold-based approaches by providing mathematically rigorous performance predictions and proactive scaling decisions.

The framework achieves significant improvements in SLA compliance, response time optimization, and oscillation elimination compared to baseline reactive scaling methods, while emphasizing operational transparency through the comprehensive real-time monitoring dashboard. Experimental evaluation demonstrated measurable enhancements in resource utilization efficiency and system stability under dynamic workload conditions characteristic of course registration systems.

The MFDS-MSO-FPR pipeline provides a mathematically grounded and practically deployable foundation for intelligent resource management in cloud-native environments, successfully bridging the gap between theoretical queuing models and operational autoscaling systems.

Continued development of such queuing theory-based systems will accelerate their transition from experimental deployments to production-grade cloud infrastructure, ultimately supporting faster, more accurate, and cost-efficient scaling decisions in latency-sensitive distributed applications while establishing a foundation for future research in intelligent autoscaling.
